\begin{frame}[fragile]\frametitle{}
\begin{center}
{\Large Prompting 101}

{\small Make everyone effective immediately}
\end{center}
\end{frame}

%%%%%%%%%%%%%%%%%%%%%%%%%%%%%%%%%%%%%%%%%%%%%%%%%%%%%%%%%%%
\begin{frame}[fragile]\frametitle{How LLMs ``Think'', Tokens and Context}
\begin{itemize}
\item LLMs process text as \textbf{tokens}, roughly 3/4 of a word on average
\item The \textbf{context window} is how much text the model can ``see'' at once, 128K to 1M tokens in modern models
\item Everything you want the model to know \textit{must be in the context}: history, instructions, documents
\item The model is \textit{stateless} by default, it forgets everything between API calls
\end{itemize}

\textbf{Practical implication:}
\begin{itemize}
\item Put the most important instructions at the \textit{beginning} and \textit{end} of the prompt (recency \& primacy effects)
\item Don't assume the model knows your codebase, provide relevant snippets
\item Token costs money: $\approx$ \$0.003 per 1K input tokens (GPT-4o pricing guide)
\end{itemize}
\end{frame}

%%%%%%%%%%%%%%%%%%%%%%%%%%%%%%%%%%%%%%%%%%%%%%%%%%%%%%%%%%%
\begin{frame}[fragile]\frametitle{Prompt Patterns, Role Prompting}
\textbf{Pattern:} Assign the model an expert role to activate relevant knowledge and style.

\begin{block}{Basic Prompt}
\texttt{Explain Kubernetes pod scheduling.}
\end{block}

\begin{block}{Role Prompt (Better)}
\texttt{You are a senior Site Reliability Engineer with 10 years of Kubernetes experience, mentoring a junior developer. Explain pod scheduling concisely, using an analogy.}
\end{block}

\textbf{Why it works:} The role activates specific patterns from training, an expert's vocabulary, precision, and communication style.\\
\textbf{Use in SDLC:} ``You are a security code reviewer following OWASP Top 10...''
\end{frame}

%%%%%%%%%%%%%%%%%%%%%%%%%%%%%%%%%%%%%%%%%%%%%%%%%%%%%%%%%%%
\begin{frame}[fragile]\frametitle{Prompt Patterns, Few-Shot Prompting}
\textbf{Pattern:} Provide 2--3 input/output examples before your actual request.

\begin{block}{Few-Shot for Commit Message Generation}
\texttt{Input: Added null check in user auth handler}\\
\texttt{Output: fix(auth): add null guard for missing user token}\\[0.4em]
\texttt{Input: Updated Dockerfile to use slim base image}\\
\texttt{Output: chore(docker): switch to python:3.11-slim base image}\\[0.4em]
\texttt{Input: Fixed off-by-one in pagination logic}\\
\texttt{Output:}
\end{block}

\textbf{Why it works:} Examples define the output format, style, and level of detail, far more precisely than any instruction.\\
\textbf{Use in SDLC:} Few-shot for test case format, release note style, JIRA ticket format.
\end{frame}

%%%%%%%%%%%%%%%%%%%%%%%%%%%%%%%%%%%%%%%%%%%%%%%%%%%%%%%%%%%
\begin{frame}[fragile]\frametitle{Prompt Patterns, Chain-of-Thought}
\textbf{Pattern:} Ask the model to \textit{reason step-by-step} before giving the final answer.

\begin{block}{Without CoT (bad for reasoning)}
\texttt{Is this SQL query safe from injection? [query]. Answer yes or no.}
\end{block}

\begin{block}{With Chain-of-Thought (better)}
\texttt{Review this SQL query for injection vulnerabilities. First, identify all user-controlled inputs. Second, check if they are parameterized. Third, check for dynamic string concatenation. Finally, give your verdict with justification.}
\end{block}

\textbf{Why it works:} Forces the model to generate intermediate reasoning, reducing errors on complex multi-step tasks. Think of it as asking someone to ``show their work.''
\end{frame}

%%%%%%%%%%%%%%%%%%%%%%%%%%%%%%%%%%%%%%%%%%%%%%%%%%%%%%%%%%%
\begin{frame}[fragile]\frametitle{Prompt Patterns, Structured Output}
\textbf{Pattern:} Ask for output in a specific structured format (JSON, YAML, Markdown table).

\begin{block}{Structured Output Prompt}
\texttt{Analyze this stack trace and return ONLY valid JSON:}\\
\texttt{\{"root\_cause": "...", "affected\_service": "...",}\\
\texttt{"suggested\_fix": "...", "severity": "P1|P2|P3"\}}
\end{block}

\textbf{Why it works for engineering:}
\begin{itemize}
\item Downstream code can parse and act on structured output programmatically
\item Reduces hallucinated free-text, enforces schema
\end{itemize}

\textbf{Advanced:} Use \textit{function calling / tool use} API feature to guarantee valid JSON schema from the model.
\end{frame}

%%%%%%%%%%%%%%%%%%%%%%%%%%%%%%%%%%%%%%%%%%%%%%%%%%%%%%%%%%%
\begin{frame}[fragile]\frametitle{Prompting for Engineering Tasks, Examples}
\begin{itemize}
\item \textbf{Debugging:}
  \texttt{``Here is a Python function and the failing test output. Identify the bug and provide a fix with explanation.''[code][output]}
\item \textbf{Refactoring:}
  \texttt{``Refactor the following function to reduce cyclomatic complexity, improve readability, and add docstrings. Preserve behavior exactly.''[code]}
\item \textbf{Writing Specs:}
  \texttt{``Given this user story, write a detailed technical specification including: API contract, data model changes, edge cases, and non-functional requirements.''}
\item \textbf{RCA Analysis:}
  \texttt{``Here are logs from the last 30 minutes before the outage. Use the 5-Whys framework to identify root cause and immediate mitigation.''[logs]}
\item \textbf{Support Investigation:}
  \texttt{``Customer reports X. Here are the relevant service logs. Identify the user-facing issue, technical cause, and ETA for fix.''[logs]}
\end{itemize}
\end{frame}


%%%%%%%%%%%%%%%%%%%%%%%%%%%%%%%%%%%%%%%%%%%%%%%%%%%%%%%%%%%
